\section{Całka McShane'a}
\providecommand{\atlasref}[2]{#2}
% Import z AM1-36, zachowane szczególne konstrukcje skryptu.
\begin{definicja}\label{it:def:mcshane}
Niech $a<b$ i niech $f:[a,b]\to\R$ ma skończone wartości. Podział znaczony $\mathcal P=\{(t_i,[x_{i-1},x_i])\}$ jest \emph{wolny $\delta$-drobny}, gdy $t_i\in[a,b]$ oraz $[x_{i-1},x_i]\subset(t_i-\delta(t_i),t_i+\delta(t_i))$, bez wymagania $t_i\in[x_{i-1},x_i]$. Funkcja $f$ jest całkowalna w sensie McShane'a, $f\in\mathcal M[a,b]$, jeśli istnieje $I\in\R$ takie, że dla każdego $\varepsilon>0$ istnieje cechowanie $\delta:[a,b]\to(0,\infty)$, dla którego każdy wolny $\delta$-drobny podział spełnia
\[
\left|S(f,\mathcal P)-I\right|<\varepsilon,\qquad
 S(f,\mathcal P)=\sum_i f(t_i)(x_i-x_{i-1}).
\]
Piszemy wówczas $I=(\mathcal M)\int_a^b f$. Na przedziale jednopunktowym przyjmujemy całkę zero.
\end{definicja}
\begin{twierdzenie}[Kryterium Cauchy'ego, liniowość i porządek]\label{it:thm:mcshane-basic}
Dla funkcji $f:[a,b]\to\R$ o skończonych wartościach
całka McShane'a istnieje wtedy i tylko wtedy, gdy dla każdego
$\varepsilon>0$ istnieje cechowanie $\delta$ takie, że
$|S(f,\mathcal P)-S(f,\mathcal Q)|<\varepsilon$ dla wszystkich
wolnych $\delta$-drobnych podziałów $\mathcal P,\mathcal Q$.
Całka jest jednoznaczna, liniowa i monotoniczna.
Na przedziale jednopunktowym przyjmujemy całkę zero.
\end{twierdzenie}
\begin{proof}
Każdy podział HK jest wolnym podziałem, zatem lemat Cousina
\ref{it:lem:kh-cousin} zapewnia istnienie wolnego podziału
dla dowolnego cechowania. Konieczność kryterium wynika
z nierówności trójkąta. Dla dostateczności wybieramy cechowania
$\delta_n$ dla błędów $2^{-n}$, zastępując je kolejno
minimami z poprzednimi. Wybieramy wolne $\delta_n$-drobne
podziały $\mathcal P_n$. Dla $m\ge n$ obie sumy są drobne
względem $\delta_n$, więc ich różnica ma moduł mniejszy
od $2^{-n}$. Sumy tworzą ciąg Cauchy'ego o granicy $I$.
Porównując dowolny wolny $\delta_n$-drobny podział
z $\mathcal P_m$ dla $m\ge n$ i przechodząc do granicy,
otrzymujemy błąd względem $I$ nie większy od $2^{-n}$. Dla dowolnej wymaganej dokładności $\varepsilon>0$ wybieramy $n$ takie, że $2^{-n}<\varepsilon$. Cechowanie $\delta_n$ spełnia wtedy definicję całki McShane'a z wartością $I$.

Jeśli dwie liczby spełniają definicję całki, wybieramy
wspólne cechowanie dla błędu $\varepsilon/2$ i jeden drobny
podział; różnica tych liczb ma moduł mniejszy od $\varepsilon$.
Dla dwóch funkcji całkowalnych $f,g$ i współczynników $\alpha,\beta\in\R$ wybieramy minimum ich cechowań dla błędu $\varepsilon/(1+|\alpha|+|\beta|)$. Liniowość skończonych sum daje
\[
 |S(\alpha f+\beta g,\mathcal P)-\alpha I_f-\beta I_g|
 \le |\alpha||S(f,\mathcal P)-I_f|
     +|\beta||S(g,\mathcal P)-I_g|<\varepsilon.
\]
To dowodzi istnienia całki kombinacji liniowej oraz jej wartości.
Dla $f\le g$ mamy $S(f,\mathcal P)\le S(g,\mathcal P)$
na każdym podziale, ponieważ długości komórek są nieujemne.
Dla $\varepsilon>0$ wybieramy wspólne cechowanie dla błędu $\varepsilon/2$ i wolny drobny podział (istniejący na mocy lematu Cousina). Wtedy $I_f\le S(f,\mathcal P)+\varepsilon/2\le S(g,\mathcal P)+\varepsilon/2\le I_g+\varepsilon$. Dowolność $\varepsilon$ daje $I_f\le I_g$.
\end{proof}

\begin{stwierdzenie}
$\mathcal R[a,b]\subset\mathcal M[a,b]\subset\HK[a,b]$, a wartości całek są zgodne.
\end{stwierdzenie}
\begin{proof}
Druga inkluzja wynika z warunków na znaczniki. Dla pierwszej ustalmy $M=\|f\|_\infty$ oraz podział $P_0$ z $U(f,P_0)-L(f,P_0)<\varepsilon/2$. Niech $B$ będzie skończonym zbiorem punktów jego podziału. Dla $t\notin B$ wybieramy $\delta(t)<\operatorname{dist}(t,B)$, zaś w punktach $t\in B$ dobieramy dodatnie cechowania tak małe, aby $4M\sum_{t\in B}\delta(t)<\varepsilon/2$. Każdy przedział wolnego podziału ze znacznikiem poza $B$ leży w tej samej komórce $P_0$ co jego znacznik. Oznaczmy komórki $P_0$ przez $J_k$, a oscylację $f$ na $J_k$ przez $\omega_k=\sup_{J_k}f-\inf_{J_k}f$. Z lokalności całki Riemanna (twierdzenie \ref{ma1:thm:riemann-additivity}) mamy dla każdej komórki $I_i$
\[
 S(f,\mathcal P)-\int_a^b f=\sum_i\left(f(t_i)|I_i|-\int_{I_i} f\right).
\]
Gdy $t_i\notin B$ oraz $t_i\in J_k$, $I_i\subset J_k$, oszacowanie całki przez kresy funkcji daje moduł składnika nie większy niż $\omega_k|I_i|$. Łączna długość tych komórek w $J_k$ jest nie większa niż $|J_k|$, więc suma błędów dla znaczników spoza $B$ jest nie większa niż $\sum_k\omega_k|J_k|=U(f,P_0)-L(f,P_0)$. Dla ustalonego $t\in B$ wszystkie komórki z tym znacznikiem mają rozłączne wnętrza i leżą w otoczeniu długości $2\delta(t)$, nawet jeśli znacznik do nich nie należy. Moduł każdego składnika jest nie większy niż $2M|I_i|$, więc ich łączny błąd nie przekracza $4M\delta(t)$. Dodając oszacowania, dostajemy błąd mniejszy niż $\varepsilon$.
\end{proof}
\noindent\emph{Porównanie w Bibliotece matematycznej:} \atlasref{AM1-36::it:thm:mcshane-absolute}{AM1-36: bezwzględna całkowalność w sensie McShane'a}.

\begin{twierdzenie}[Bezwzględna całkowalność]\label{it:thm:mcshane-absolute}
Jeżeli $f\in\mathcal M[a,b]$, to $|f|\in\mathcal M[a,b]$. Tak więc całka McShane'a jest bezwzględna, mimo że definiujemy ją przez sumy znaczonych przedziałów.
\end{twierdzenie}
\begin{proof}
Niech $I=\int_a^b f$ w sensie McShane'a. Wybieramy cechowanie $\delta$ zapewniające błąd mniejszy niż $\varepsilon/3$. Dla dwóch dowolnych wolnych $\delta$-drobnych podziałów $\mathcal P=\{(t_i,I_i)\}$ i $\mathcal Q=\{(s_j,J_j)\}$ rozważamy wspólne zagęszczenie $K_{ij}=I_i\cap J_j$, pomijając przecięcia zerowej długości. Każda komórka $K_{ij}$ leży w otoczeniach cechowania zarówno $t_i$, jak i $s_j$, więc obydwa znaczniki są dopuszczalne. Wybierając w jednej sumie znacznik z większą wartością $f$, a w drugiej z mniejszą, otrzymujemy dwie dopuszczalne sumy $S_+,S_-$, zatem
\[
\sum_{i,j}|f(t_i)-f(s_j)|\,|K_{ij}|=S_+-S_-<2\varepsilon/3.
\]
Ponieważ $|I_i|=\sum_j|K_{ij}|$, $|J_j|=\sum_i|K_{ij}|$ i $\bigl||u|-|v|\bigr|\le|u-v|$, mamy
$|S(|f|,\mathcal P)-S(|f|,\mathcal Q)|\le\sum_{i,j}|f(t_i)-f(s_j)|\,|K_{ij}|<\varepsilon$.
Kryterium Cauchy'ego z twierdzenia \ref{it:thm:mcshane-basic} dowodzi więc istnienia całki McShane'a z $|f|$. Dla HK podobny argument nie działa: znacznik pierwotnej komórki nie musi należeć do jej przecięcia z inną komórką, co narusza warunek HK.
\end{proof}

\begin{przyklad}[Inkluzje są właściwe]
Funkcja Dirichleta należy do $\mathcal M[0,1]$ (lemat \ref{ma1:kh:lem:kh-null-modification} działa także dla wolnych znaczników), ale nie do $\mathcal R[0,1]$, ponieważ jest nieciągła w każdym punkcie (twierdzenie \ref{thm:riemann-criterion}). Funkcja $x^{-1/2}$ po określeniu w $0$ również należy do $\mathcal M\setminus\mathcal R$. Pochodna niecałkowalna bezwzględnie z rozdziału o całce HK należy do $\HK\setminus\mathcal M$; wynika to z twierdzenia \ref{it:thm:mcshane-lebesgue}. To samo twierdzenie uzasadnia całkowalność McShane'a funkcji $x^{-1/2}$, której całka Lebesgue'a wynosi $2$: funkcje $x^{-1/2}\chi_{[1/n,1]}(x)$, przy wartości zero w $0$, rosną prawie wszędzie do tej funkcji i mają całki $2(1-1/\sqrt n)$, co przez twierdzenie \ref{thm:mct-ae} daje żądaną wartość. Brak całkowalności Riemanna wynika z nieograniczoności.
\end{przyklad}

\begin{uwaga}[Kolejność rozszerzeń]
Całka McShane'a rozszerza Riemanna i jest zawarta w HK. W teorii funkcji rzeczywistych opisuje dokładnie całkę Lebesgue'a; wykażemy to dalej. Wolność znaczników powiększa rodzinę sum, wobec których wymagamy jednakowych oszacowań, więc \emph{zmniejsza} klasę funkcji całkowalnych w porównaniu z HK.
\end{uwaga}

% ============================================================
% UZUPELNIENIE MATEMATYCZNE: ROZDZIAL 7
% ============================================================
\subsection*{Konstrukcja całki McShane'a dla funkcji charakterystycznych}
\begin{twierdzenie}[Funkcje charakterystyczne zbiorów mierzalnych]\label{tw:mcshane-indykator}
Niech $E\subset[a,b]$ będzie zbiorem mierzalnym w sensie Lebesgue'a. Wówczas $\ind E$ jest całkowalna w sensie McShane'a i
$(\mathcal M)\int_a^b\chi_E(t)\,\mathrm dt=\lambda_1(E)$.
\end{twierdzenie}
\begin{proof}
Ustalmy $\varepsilon>0$. Z regularności miary Lebesgue'a wybieramy zbiór zwarty $K\subset E$ i zbiór otwarty $O\supset E$ (w topologii względnej $[a,b]$), dla których $\lambda_1(O\setminus K)<\varepsilon$. Konstruujemy cechowanie: dla $t\in E$ wybieramy $\delta(t)>0$ tak, by $(t-\delta(t),t+\delta(t))\cap[a,b]\subset O$; dla $t\notin E$ dobieramy je tak, by wymienione otoczenie nie przecinało $K$ (jest to możliwe, ponieważ $K$ jest domknięty i $t\notin K$).

Niech $\mathcal P$ będzie dowolnym wolnym $\delta$-drobnym podziałem. Każda jego komórka z oznaczeniem $t\in E$ leży w $O$, a każda komórka ze znacznikiem $t\notin E$ nie przecina $K$. Ponieważ komórki pokrywają cały przedział, suma długości komórek ze znacznikami w $E$ mieści się między $\lambda_1(K)$ a $\lambda_1(O)$. Jest ona dokładnie sumą McShane'a $S(\ind E,\mathcal P)$. Wobec $\lambda_1(K)\le\lambda_1(E)\le\lambda_1(O)$ otrzymujemy $|S-\lambda_1(E)|<\varepsilon$.
\end{proof}
\begin{wniosek}[Funkcje proste]\label{wn:mcshane-proste}
Niech $s=\sum_{k=1}^m c_k\ind{E_k}$, gdzie $c_k\in\R$ i wszystkie $E_k\subset[a,b]$ są mierzalne. Wtedy $s\in\mathcal M[a,b]$ oraz $(\mathcal M)\int_a^b s(t)\,\mathrm dt=\sum_{k=1}^m c_k\lambda_1(E_k)$.
\end{wniosek}
\begin{proof}
Każdy składnik jest całkowalny na mocy twierdzenia \ref{tw:mcshane-indykator}; liniowość sum jest bezpośrednią konsekwencją definicji całki McShane'a.
\end{proof}
\begin{przyklad}[Zbiór Cantora]
Dla klasycznego zbioru Cantora $K\subset[0,1]$ mamy $\lambda_1(K)=0$, więc $(\mathcal M)\int_0^1\chi_K(t)\,\mathrm dt=0$. Ten sam wynik zachodzi dla funkcji Dirichleta z racji $\lambda_1(\Q\cap[0,1])=0$. Funkcja charakterystyczna klasycznego zbioru Cantora jest też całkowalna w sensie Riemanna, ponieważ zbiór jej nieciągłości jest zawarty w zbiorze Cantora i ma miarę zero. Funkcja Dirichleta nie jest całkowalna w sensie Riemanna, gdyż jest nieciągła w każdym punkcie.
\end{przyklad}
\begin{uwaga}[Zakres uogólnienia]
Twierdzenie o funkcjach charakterystycznych daje od razu zgodność całki McShane'a z całką Lebesgue'a na funkcjach prostych. Przejście do wszystkich funkcji z $L^1$ nie jest wyłącznie konsekwencją liniowości: wymaga dodatkowej kontroli sum dla przybliżeń i jest przeprowadzone w twierdzeniu \ref{it:thm:mcshane-lebesgue} w niniejszym rozdziale.
\end{uwaga}

\subsection{Własności i twierdzenia graniczne}
\begin{wniosek}[Działania kratowe]\label{it:cor:mcshane-lattice}
Dla $f,g\in\mathcal M[a,b]$ także $f^+$, $f^-$,
$\max(f,g)$ i $\min(f,g)$ należą do $\mathcal M[a,b]$.
\end{wniosek}
\begin{proof}
Stosujemy twierdzenia \ref{it:thm:mcshane-basic}
i \ref{it:thm:mcshane-absolute} do tożsamości
$f^+=(|f|+f)/2$ i $f^-=(|f|-f)/2$
oraz $\max(f,g)=(f+g+|f-g|)/2$ i
$\min(f,g)=(f+g-|f-g|)/2$.
\end{proof}

\begin{twierdzenie}[Wspólne cechowanie]\label{it:thm:mcshane-equiintegrability}
Niech $f_n\in\mathcal M[a,b]$ z całkami $I_n$ i niech
$f_n(t)\to f(t)\in\R$ dla każdego $t$.
Załóżmy, że dla każdego $\varepsilon>0$ istnieje jedno
cechowanie $\delta$ takie, że $|S(f_n,\mathcal P)-I_n|<\varepsilon$
dla wszystkich $n$ i wszystkich wolnych $\delta$-drobnych
podziałów $\mathcal P$. Wówczas $f\in\mathcal M[a,b]$
oraz $I_n\to(\mathcal M)\int_a^b f$.
\end{twierdzenie}
\begin{proof}
Dla ustalonego $\varepsilon>0$ wybieramy wspólne cechowanie
dla błędu $\varepsilon/3$ i jeden wolny drobny podział.
Jego skończona suma dla $f_n$ zbiega do sumy dla $f$,
ponieważ zbieżność zachodzi we wszystkich znacznikach.
Dla dostatecznie dużych $n,m$ moduł różnicy sum jest
mniejszy od $\varepsilon/3$, więc z nierówności trójkąta
$|I_n-I_m|<\varepsilon$. Niech $I=\lim_n I_n$.
Dla dowolnej dokładności $\eta>0$ wybieramy wspólne
cechowanie dla błędu $\eta/2$. Dla każdego wolnego
drobnego podziału $\mathcal Q$ przejście do granicy
w nierówności $|S(f_n,\mathcal Q)-I_n|<\eta/2$ daje
$|S(f,\mathcal Q)-I|\le\eta/2<\eta$.
\end{proof}

\begin{wniosek}[Zbieżność jednostajna]\label{it:cor:mcshane-uniform}
Jeżeli $f_n\in\mathcal M[a,b]$ i $f_n\to f$ jednostajnie,
to $f\in\mathcal M[a,b]$ i całki zbiegają.
\end{wniosek}
\begin{proof}
Dla dostatecznie dużych $n,m$ różnica $f_n-f_m$ jest
ograniczona. Z oszacowania sum oraz liniowości wynika
$|I_n-I_m|\le(b-a)\|f_n-f_m\|_\infty$.
Stąd $I=\lim I_n$ istnieje. Dla $\varepsilon>0$ wybieramy
$n$ tak, aby $(b-a)\|f-f_n\|_\infty<\varepsilon/3$
oraz $|I_n-I|<\varepsilon/3$. Na wolnych podziałach
drobnych względem cechowania dla $f_n$ z błędem
$\varepsilon/3$ mamy
\[
 |S(f,\mathcal P)-I|
 \le (b-a)\|f-f_n\|_\infty+|S(f_n,\mathcal P)-I_n|
     +|I_n-I|<\varepsilon.
\]
Nie zakładamy z góry całkowalności granicy ani ograniczoności
poszczególnych funkcji.
\end{proof}

\begin{wniosek}[Lokalność, zbieżność monotoniczna i zmajoryzowana]
\label{it:cor:mcshane-lebesgue-consequences}
Całkowalność McShane'a na $[a,b]$ pociąga całkowalność
na każdym podprzedziale; całka jest addytywna przy
sklejaniu dwóch sąsiednich przedziałów.
Jeśli $f_n\in\mathcal M[a,b]$ oraz $f_n\to f$ prawie
wszędzie, zbieżność całek i całkowalność $f$ zachodzą
w każdym z następujących przypadków:
\begin{enumerate}[label=(\alph*)]
\item $|f_n|\le g$ prawie wszędzie dla pewnej
$g\in\mathcal M[a,b]$;
\item $f_n\uparrow f$ prawie wszędzie i $\sup_n I_n<\infty$,
gdzie $I_n=(\mathcal M)\int_a^b f_n$.
\end{enumerate}
Zakładamy skończone wartości funkcji $f$; na zbiorze
wyjątkowym wolno nadać jej dowolne skończone wartości.
\end{wniosek}
\begin{proof}
Najpierw dowodzimy lokalności bez reprezentacji miarowej. Dla $a<c<b$ wybieramy cechowanie z kryterium Cauchy'ego na $[a,b]$. Dwa wolne drobne podziały $[a,c]$ uzupełniamy tym samym drobnym podziałem $[c,b]$, istniejącym dzięki lematowi Cousina. Po odjęciu wspólnej sumy kryterium Cauchy'ego daje całkowalność na $[a,c]$. Zamiana ról i powtórzenie argumentu daje całkowalność na każdym podprzedziale. Odwrotnie, jeżeli istnieją całki na obu częściach, wybieramy na nich cechowania dla błędu $\varepsilon/2$. Dla $t\ne c$ ograniczamy odpowiednie cechowanie przez $|t-c|/2$, a dla $t=c$ bierzemy minimum obu cechowań. Żadna komórka ze znacznikiem $t\ne c$ nie może leżeć po przeciwnej stronie $c$ niż jej znacznik, nawet gdy znacznik nie należy do komórki. Komórkę przecinającą $c$ rozcinamy w $c$, zachowując znacznik $c$. Po rozcięciu otrzymujemy wolne drobne podziały obu części o niezmienionej łącznej sumie. Nierówność trójkąta daje błąd mniejszy od $\varepsilon$ względem sumy całek lokalnych. To dowodzi również addytywności.

Dla twierdzeń granicznych stosujemy identyfikację \ref{it:thm:mcshane-lebesgue}: wszystkie funkcje McShane'a są całkowalne w sensie Lebesgue'a i odwrotnie. Wybieramy jeden mierzalny zbiór miary zero zawierający wszystkie zbiory wyjątkowe z założeń. Poza nim $f$ jest granicą funkcji mierzalnych; z zupełności miary Lebesgue'a wynika mierzalność $f$. W (a) stosujemy twierdzenie \ref{thm:dct} z majorantą $|g|$. W (b) funkcje $f_n-f_1$ są nieujemne prawie wszędzie i rosną do $f-f_1$ prawie wszędzie. Ich całki są ograniczone przez $\sup_n I_n-I_1$, więc twierdzenie \ref{thm:mct-ae} daje całkowalność $f-f_1$ i zbieżność jego całek. Dodając $f_1$, dostajemy całkowalność $f$ i zbieżność $I_n$.
Ponowna identyfikacja całek przez twierdzenie
\ref{it:thm:mcshane-lebesgue} kończy dowód.
\end{proof}

\subsection{Zgodność całek McShane'a, Lebesgue'a i HK}
\begin{lemat}[Bezwzględna ciągłość całki]\label{it:lem:integral-absolute-continuity}\label{ma1:kh:lem:integral-ac}
Jeśli $f\in\mathcal L^1(\mu)$, to dla każdego $\varepsilon>0$ istnieje $\eta>0$ takie, że $\mu(A)<\eta$ pociąga $\int_A|f|\,\mathrm d\mu<\varepsilon$ dla każdego $A\in\mathcal F$ w przestrzeni z miarą $(X,\mathcal F,\mu)$.
\end{lemat}
\begin{proof}
Wybieramy $M>0$ z $\int_{\{|f|>M\}}|f|\,\mathrm d\mu<\varepsilon/2$ oraz $\eta=\varepsilon/(2M)$. Wtedy $\int_A|f|\,\mathrm d\mu\le M\mu(A)+\int_{\{|f|>M\}}|f|\,\mathrm d\mu<\varepsilon$. Istotnie, funkcje $|f|\chi_{\{|f|\le n\}}$ rosną prawie wszędzie do $|f|$ (wartości nieskończone występują tylko na zbiorze zerowym). Z twierdzenia \ref{thm:mct-ae} ich całki dążą do skończonej całki $|f|$. Z addytywności całki po rozłącznych zbiorach otrzymujemy $\int_{\{|f|>n\}}|f|\,\mathrm d\mu\to0$, co uzasadnia wybór $M$.
\end{proof}

\begin{twierdzenie}[McShane = Lebesgue]\label{it:thm:mcshane-lebesgue}
Niech $a<b$. Dla funkcji rzeczywistych zachodzi
$\mathcal M[a,b]=\mathcal L^1([a,b],\lambda_1)$
jako równość zbiorów konkretnych funkcji o skończonych
wartościach, gdzie $\lambda_1$ jest miarą uzupełnioną.
Dowolne zmiany skończonych wartości na zbiorach
miary zero zachowują całkowalność i wartość całki.
Wartości obu całek są równe.
\end{twierdzenie}
\begin{proof}
\emph{Krok 1: Lebesgue $\Rightarrow$ McShane.} Niech $f\in\mathcal L^1$ i $\varepsilon>0$. Wybieramy $h>0$ z $h(b-a)<\varepsilon/3$ i rozbijamy dziedzinę na mierzalne zbiory $E_k=\{kh\le f<(k+1)h\}$, $k\in\mathbb Z$. Z bezwzględnej ciągłości całki wybieramy $\eta>0$ tak, że $\lambda_1(A)<\eta$ pociąga $\int_A|f|\,\mathrm d\lambda_1<\varepsilon/3$. Ustalamy dodatnie liczby $w_k=2^{-|k|}/3$ ($k\in\Z$), których suma wynosi $1$. Z regularności miary wybieramy względnie otwarte $G_k\supset E_k$ z
\[
 \lambda_1(G_k\setminus E_k)<
 w_k\min\left\{\eta/2,\frac{\varepsilon}{6(1+|kh|+h)}\right\}.
\]
Dla $U=\bigcup_k(G_k\setminus E_k)$ mamy wtedy $\lambda_1(U)\le\eta/2$ i $\int_U|f|\,\mathrm d\lambda_1<\varepsilon/3$, a suma $\sum_k(|kh|+h)\lambda_1(G_k\setminus E_k)$ nie przekracza $\varepsilon/6$. Dla $t\in E_k$ wybieramy $\delta(t)>0$ z $(t-\delta(t),t+\delta(t))\cap[a,b]\subset G_k$.

Niech $\mathcal P=\{(t_i,I_i)\}$ będzie dowolnym wolnym $\delta$-drobnym podziałem. Zatem $I_i\subset G_k$ dla $t_i\in E_k$, również gdy $t_i\notin I_i$. Ponieważ komórki mają rozłączne wnętrza, a ich wspólne końce są zbiorami zerowymi, z addytywności całki Lebesgue'a mamy
\[
 S(f,\mathcal P)-\int_a^b f(x)\,\mathrm d\lambda_1(x)
 =\sum_i\int_{I_i}\bigl(f(t_i)-f(x)\bigr)\,\mathrm d\lambda_1(x).
\]
Przy znaczniku $t_i\in E_k$ na części $I_i\cap E_k$ moduł różnicy wartości $f(t_i)$ i $f(x)$ jest mniejszy od $h$, ponieważ obie wartości należą do $[kh,(k+1)h)$. Pozostałe mierzalne części $I_i\setminus E_k$ leżą w $U$; ich przecięcia zawierają co najwyżej wspólne końce komórek, a więc mają miarę zero. Dla ustalonego $k$ suma ich miar po komórkach ze znacznikami w $E_k$ nie przekracza $\lambda_1(G_k\setminus E_k)$. Po rozdzieleniu sumy otrzymujemy
\[
\left|S(f,\mathcal P)-\int_a^b f\,\mathrm d\lambda_1\right|
\le h(b-a)+\sum_k(|kh|+h)\lambda_1(G_k\setminus E_k)+\int_U|f|\,\mathrm d\lambda_1<\varepsilon.
\]

\emph{Krok 2: McShane $\Rightarrow$ Lebesgue.} Ponieważ każdy podział ze znacznikami należącymi do komórek jest również podziałem z wolnymi znacznikami, całkowalność McShane'a pociąga całkowalność HK z tą samą wartością całki. Mierzalność $f$ wynika z twierdzenia \ref{it:thm:kh-primitive-derivative}. Z twierdzenia \ref{it:thm:mcshane-absolute} także $|f|\in\mathcal M$, z całką $J\ge0$. Obcięcia $h_n=\min(|f|,n)$ są ograniczone i mierzalne, zatem Lebesgue'a całkowalne, a na mocy kroku 1 także McShane'a całkowalne. Monotoniczność całki McShane'a (twierdzenie \ref{it:thm:mcshane-basic}) i zgodność z kroku 1 dają $\int h_n\,\mathrm d\lambda_1\le J$. Przez zbieżność monotoniczną $\int|f|\,\mathrm d\lambda_1\le J<\infty$. Krok 1 i jednoznaczność całki McShane'a
(twierdzenie \ref{it:thm:mcshane-basic})
dowodzą zgodności wartości. Jeżeli zmienimy
skończone wartości $f$ na zbiorze miary zero,
to zupełność miary zapewnia mierzalność zmodyfikowanej
funkcji, a twierdzenie \atlasref{ME-07::thm:ae-integral-invariance}{ME-07: niezmienniczość całki przy równości prawie wszędzie}
zapewnia równość całek Lebesgue'a i zachowanie
całkowalności. Z wykazanej równoważności wynika to
samo dla McShane'a. Ponieważ każdy podział HK jest
wolnym podziałem, w szczególności dla każdej
$f\in\mathcal L^1$ zachodzi $f\in\HK$ i
$(\HK)\int_a^b f=\int_{[a,b]}f\,\mathrm d\lambda_1$.
Nie korzystaliśmy tutaj z kryterium bezwzględnej
całkowalności HK, które udowodnimy poniżej.
Dla $a=b$ każda skończona funkcja na singletonie
jest mierzalna, całkowalna i obie całki wynoszą zero.
\end{proof}
\begin{twierdzenie}[Kryterium bezwzględnej całkowalności HK]\label{ma1:kh:thm:kh-absolute}\label{it:thm:hk-absolute-integrability}
Niech $f:[a,b]\to\R$ będzie funkcją o skończonych
wartościach na przedziale zwartym. Jeżeli $f\ge0$
jest HK-całkowalna, to $f\in\mathcal L^1$.
Dla dowolnej takiej funkcji mamy
\[
f\in\mathcal L^1\quad\Longleftrightarrow\quad f\in\HK\ \text{i}\ |f|\in\HK.
\]
Każda ograniczona funkcja na $[a,b]$ jest
HK-całkowalna wtedy i tylko wtedy, gdy jest
mierzalna w sensie Lebesgue'a. W przypadku Riemanna
konieczny i wystarczający jest warunek, aby zbiór
punktów nieciągłości miał miarę zero
(twierdzenie \ref{thm:riemann-criterion}).
\end{twierdzenie}
\begin{proof}
Z twierdzenia \ref{it:thm:kh-primitive-derivative} funkcja HK-całkowalna jest mierzalna. Dla nieujemnej $f$ obcięcia $h_n=\min(f,n)$ są mierzalne, ograniczone i całkowalne Lebesgue'a, więc należą także do HK. Z monotoniczności HK mamy $\int h_n\,\mathrm d\lambda_1\le(\HK)\int f<\infty$. Zbieżność monotoniczna daje $f\in\mathcal L^1$. Zastosowanie pierwszej części do $|f|$ oraz inkluzja $\mathcal L^1\subset\HK$ dają równoważność. Ostatnia teza wynika z mierzalności funkcji HK i skończoności miary przedziału.
\end{proof}

\begin{wniosek}[Części dodatnia i ujemna w kryterium HK]
\label{it:cor:hk-positive-negative-integrability}
Dla każdej funkcji rzeczywistej $f$ o skończonych
wartościach na $[a,b]$ następujące warunki są równoważne:
\begin{enumerate}[label=(\roman*)]
\item $f\in\mathcal L^1([a,b],\lambda_1)$;
\item $f^+,f^-\in\HK[a,b]$;
\item $f,|f|\in\HK[a,b]$;
\item $f\in\mathcal M[a,b]$.
\end{enumerate}
Jeżeli zachodzi którykolwiek z nich, to
\[
 (\HK)\int_a^b|f|=
 \int_{[a,b]}|f|\,\mathrm d\lambda_1
 =(\HK)\int_a^b f^++(\HK)\int_a^b f^-.
\]
Jeśli natomiast $f\in\HK[a,b]\setminus\mathcal L^1$,
to żadna z funkcji $f^+,f^-$ nie jest HK-całkowalna,
a obie ich całki Lebesgue'a wynoszą $+\infty$.
\end{wniosek}
\begin{proof}
Z (i) wynika (ii), ponieważ $0\le f^\pm\le|f|$,
więc obie części są Lebesgue'a całkowalne, a
$\mathcal L^1\subset\HK$ z twierdzenia
\ref{it:thm:mcshane-lebesgue} i jego dowodu.
Z (ii) wynika (iii) przez liniowość HK i tożsamości
$f=f^+-f^-$, $|f|=f^++f^-$. Implikacja
(iii)$\Rightarrow$(i) jest treścią twierdzenia
\ref{it:thm:hk-absolute-integrability}, a równoważność
(i) i (iv) wynika z twierdzenia
\ref{it:thm:mcshane-lebesgue}. Wzór na całkę modułu
otrzymujemy z zgodności wartości całek i addytywności
całki nieujemnej.

Jeżeli $f\in\HK$ i jedna z funkcji $f^+,f^-$
należy do HK, to druga również należy do HK,
ponieważ $f^-=f^+-f$, $f^+=f+f^-$.
Z (ii)$\Rightarrow$(i) wynika wówczas $f\in\mathcal L^1$.
Dla $f\in\HK\setminus\mathcal L^1$ obie części są
mierzalne (twierdzenie \ref{it:thm:kh-primitive-derivative}),
lecz żadna nie może mieć skończonej całki Lebesgue'a,
gdyż wówczas należałaby do HK. Stąd obie całki
Lebesgue'a są nieskończone. Całka Lebesgue'a ze
znakiem nie jest wtedy określona, mimo istnienia
skończonej całki HK.
\end{proof}

